% docs/physics/latex/template_reaction.tex - TEMPLATE of a SIMC++ physics note for a reaction or physics-model plug-in.
%
% How to use it (docs/physics/README.md, "Writing a note"):
%   1. copy this file to docs/physics/latex/NN_<name>.tex (NN = the next free number);
%   2. add the note to docs/physics/latex/notes.json: {"source", "output": "SIMCpp_Physics_NN_<name>", "title",
%      "copies": [the plug-in folders: "src/reactions/<r>" if it exists, "src/physics/<r>"]};
%   3. add or update the README.md of every copy folder (it must name the PDF; test docs_physics_pdf_copies);
%   4. replace every <...> below, keep the nine sections in this order, delete these instructions;
%   5. build: python3 -I docs/physics/build_physics_pdf.py --source-root . --build-dir /tmp/p --note NN_<name>.tex,
%      then cmake --build <build> --target physics_pdf to publish the PDF, its copies and the book.
%
% Rules (enforced by build_physics_pdf.py unless marked "review"):
%   * every display: \begin{feq}{<C++ file: function>}{<Fortran file:lines>} ... \end{feq}, or feqa for an align
%     block; C++ first, Fortran second; write "none" as the Fortran citation of a SIMC++-only formula;
%   * labels with \lbl{eq:...}, \lbl{sec:...}, \lbl{tab:...}; references \eqr, \secr, \tabr; to note 00 \ceqr, \csecr,
%     \ctabr; no \label;
%   * literature in the simcrefs environment, cited with \citer{key};
%   * finding ids of FORTRAN_DIFFERENCES.md only inside the simcfindings table of section 9;
%   * (review) tables: ONE ITEM PER ROW (one input, one column, one routine, one finding, one test per row);
%   * (review) physics notation in the text (E'_min, not edge%e%E%min); use the symbols and terms of note 00
%     (tables "Symbols" and "Glossary"); define any new symbol where it first appears;
%   * (review) the reaction prime is \pr: H(e,e\pr p);
%   * (review) the text describes SIMC++ first; the Fortran appears in the citations and in section 9.
\documentclass{simcnote}
\simccommon            % the labels of the common note 00 (needed for \ceqr, \csecr, \ctabr)
\begin{document}
\simcnotemeta{NN}{<short title for the header>}{<Title with \pr where a prime is needed>}%
  {Reaction plug-in \texttt{<registered name>} with the physics model \texttt{<registered name>}: <one line>}
\simcnotemaketitle
\simcabstract{<Three to five sentences: the process, what is generated, the cross-section model, what is special.>}
\simcnotetoc

% =====================================================================================================================
\section{The reaction in SIMC++}\lbl{sec:reaction}
\subsection{Process and scope}
% The process; what is simulated; what is not simulated (approximations, missing channels).
<text>

\subsection{Components}
% One row per component: reaction (IReaction), base class if any, physics model (IPhysicsModel), kernels, spectral
% function or data tables, and the shared stages (refer to \ctabr{tab:plugins}).
\begin{simctable}{footnotesize}{L{0.16\linewidth}L{0.17\linewidth}L{0.17\linewidth}L{0.24\linewidth}L{0.26\linewidth}}
\caption{The SIMC++ components of <reaction> (one component per row).}\lbl{tab:components}\\
\toprule\rowcolor{simchead}\simcth{role} & \simcth{interface} & \simcth{registered name} & \simcth{class} & \simcth{file}\\\midrule\endfirsthead
\toprule\rowcolor{simchead}\simcth{role} & \simcth{interface} & \simcth{registered name} & \simcth{class} & \simcth{file}\\\midrule\endhead
\bottomrule\endlastfoot
reaction & \code{IReaction} & \code{<name>} & \code{<Class>} & \file{reactions/<dir>/<File>.cpp}\\
cross section & \code{IPhysicsModel} & \code{<name>} & \code{<ns::Class>} & \file{physics/<dir>/<File>.cpp}\\
\end{simctable}

\subsection{Selection and override}
% Which input flags select the plug-ins (\ctabr{tab:defaults}), the masses and flags the set-up derives, and the JSON
% keys of the reaction and the model (a table, one key per row) with an example.
<text>

\subsection{Event flow}
% One row per step of \ctabr{tab:chain}: what is specific to this reaction.
\begin{simctable}{footnotesize}{L{0.05\linewidth}L{0.33\linewidth}L{0.62\linewidth}}
\caption{One try of <reaction> (one step per row).}\lbl{tab:flow}\\
\toprule\rowcolor{simchead}\simcth{\#} & \simcth{C++ call} & \simcth{what is specific to this reaction}\\\midrule\endfirsthead
\bottomrule\endlastfoot
1 & \code{StandardBeam::generate} & <...>\\
\end{simctable}

% =====================================================================================================================
\section{Kinematics, generated variables and generation volume}\lbl{sec:gen}
% The generated variables (a table, one variable per row), their ranges and limits (equations), the generation weight
% and the generation volume.
\begin{feq}{reactions/<dir>/<File>.cpp: <Class>::generate}{event.f:<lines>}
<\text{formula}>
\lbl{eq:genvol}
\end{feq}

% =====================================================================================================================
\section{Completing the event}\lbl{sec:complete}
\subsection{Vertex}
% How the unmeasured quantities are solved for (equations, roots, failure conditions).
\subsection{Reconstruction}
% The reconstructed quantities specific to the reaction (E_m, t, ...).

% =====================================================================================================================
\section{Jacobians}\lbl{sec:jac}
% Every factor of J, with its derivation when it is not obvious.

% =====================================================================================================================
\section{Cross-section model}\lbl{sec:xsec}
\subsection{What the reaction requires from a model}
% The IPhysicsModel methods used, the vertex quantities read, the units of the returned cross section, and what a
% replacement model must provide.
\subsection{The model \code{<name>}}
% ALL formulas evaluated, with parameter tables (one parameter per row), data files and interpolations; the literature
% the code cites (\citer), and where the code differs from it.

% =====================================================================================================================
\section{Radiation, energy loss and decays}\lbl{sec:rad}
% Only what is specific to the reaction: photon limits, the hadron's radiation, decays in flight.

% =====================================================================================================================
\section{Inputs and outputs}\lbl{sec:io}
\subsection{Inputs}
\begin{simctable}{footnotesize}{L{0.24\linewidth}L{0.12\linewidth}L{0.08\linewidth}L{0.56\linewidth}}
\caption{Inputs specific to <reaction> (one input per row).}\lbl{tab:inputs}\\
\toprule\rowcolor{simchead}\simcth{input} & \simcth{unit} & \simcth{default} & \simcth{meaning here}\\\midrule\endfirsthead
\bottomrule\endlastfoot
\code{<input>} & <unit> & <default> & <meaning>\\
\end{simctable}
\subsection{Outputs}
% The tree layout: one column per row (# | name | unit | content); reuse a frag_*.tex if the layout is shared.
\begin{simctable}{footnotesize}{L{0.06\linewidth}L{0.15\linewidth}L{0.12\linewidth}L{0.67\linewidth}}
\caption{Tree columns of <reaction> (one column per row).}\lbl{tab:columns}\\
\toprule\rowcolor{simchead}\simcth{\#} & \simcth{name} & \simcth{unit} & \simcth{content}\\\midrule\endfirsthead
\bottomrule\endlastfoot
1 & \code{<name>} & <unit> & <content>\\
\end{simctable}

% =====================================================================================================================
\section{Validation in SIMC++}\lbl{sec:validation}
% The CTest tests that cover the reaction (one test per row: name | labels | what it checks); numbers only when they can
% be cited from the validation ledger or a test.
\begin{simctable}{footnotesize}{L{0.44\linewidth}L{0.17\linewidth}L{0.39\linewidth}}
\caption{CTest tests that cover <reaction> (one test per row).}\lbl{tab:tests}\\
\toprule\rowcolor{simchead}\simcth{test} & \simcth{labels} & \simcth{what it checks}\\\midrule\endfirsthead
\bottomrule\endlastfoot
\code{<test>} & <labels> & <what>\\
\end{simctable}

% =====================================================================================================================
\section{Fortran origin and Fortran findings}\lbl{sec:origin}
\begin{simctable}{footnotesize}{L{0.42\linewidth}L{0.58\linewidth}}
\caption{Fortran routine $\to$ C++ function (one routine per row).}\lbl{tab:map}\\
\toprule\rowcolor{simchead}\simcth{Fortran (upstream 59ccdf6)} & \simcth{C++ (src/...)}\\\midrule\endfirsthead
\bottomrule\endlastfoot
<file.f routine> & \csrc{<dir/File.cpp: Function>}\\
\end{simctable}

\begin{simcfindings}{footnotesize}{L{0.18\linewidth}L{0.5\linewidth}L{0.32\linewidth}}
\caption{Fortran findings that touch <reaction> (one finding per row).}\lbl{tab:findings}\\
\toprule\rowcolor{simchead}\simcth{id} & \simcth{Fortran behaviour} & \simcth{default / compat}\\\midrule\endfirsthead
\bottomrule\endlastfoot
\finding{<id>} & <behaviour> & <default> / <compat>\\
\end{simcfindings}

\begin{simcrefs}
\simcref{<key>} <reference> (cited at \fsrc{<file.f:line>}).
\end{simcrefs}
\end{document}
